% 水星（综述）
% license CCBYSA3
% type Wiki

本文根据 CC-BY-SA 协议转载翻译自维基百科\href{https://en.wikipedia.org/wiki/Mercury_(planet)}{相关文章}。

\begin{figure}[ht]
\centering
\includegraphics[width=8cm]{./figures/8ad510c3701a820c.png}
\caption{} \label{fig_Mercur_1}
\end{figure}
Mercury 是距离太阳最近的行星，也是太阳系中最小的行星。它是一颗岩质行星，具有极其稀薄的大气层，其表面重力略高于火星。水星的表面与地球的月球相似，布满大量撞击坑，并具有由逆断层形成的广阔断崖系统（rupes），以及由抛射物形成的明亮光条结构。其最大撞击盆地——卡洛里斯平原（Caloris Planitia）——直径为 \(1{,}550\,\mathrm{km}\)（\(960\,\mathrm{mi}\)），约为整颗行星直径 \(4{,}880\,\mathrm{km}\)（\(3{,}030\,\mathrm{mi}\)）的三分之一。作为最靠内侧运行的行星，它在地球的天空中总是接近太阳出现，呈现为“晨星”或“昏星”。它也是从地球前往所需 \(\Delta v\) 最大的行星，同时亦是往返太阳系中其他行星时所需能量最高者。

水星的恒星年（\(88.0\) 个地球日）与恒星日（\(58.65\) 个地球日）构成 \(3:2\) 的自转—公转共振。因此，一个太阳日（从日出到下一次日出）约持续 \(176\) 个地球日，即其恒星年的两倍。这意味着水星某一侧将在一个水星年（\(88\) 个地球日）内持续处于日照之下；而在接下来的轨道周期中的另一个 \(88\) 个地球日内，该侧将一直处于黑暗之中，直到下一次日出。行星表面之上存在极度稀薄的外逸层以及微弱但足以偏转太阳风的磁场。由于水星轨道偏心率极高，其表面所接受的日照强度与温度变化极为剧烈：赤道地区夜间温度可达 \(-170^{\circ}\mathrm{C}\)（\(-270^{\circ}\mathrm{F}\)），而白昼可升至 \(420^{\circ}\mathrm{C}\)（\(790^{\circ}\mathrm{F}\)）。由于轴倾角极小，水星两极区域永久处于阴影之中，这强烈暗示其撞击坑内可能存在水冰。

与太阳系其他行星相同，水星形成于约 \(4.5\times10^{9}\) 年前。关于其起源与演化存在诸多相互竞争的假说，其中一些涉及与微行星碰撞及岩石蒸发过程；截至 2020 年代初，水星地质历史的许多宏观细节仍在研究中，或需等待更多深空探测数据。其地幔成分高度均一，这表明水星早期可能存在类似月球的岩浆海。根据当前模型，水星可能具有固态硅酸盐地壳与地幔，其下为固态外核、更深处的液态内核层，以及固态内核。约在七至八十亿年后，随着太阳演化为红巨星，水星预计将被毁灭，金星亦然，地球与月球也可能遭遇同样命运\(^\text{[20]}\)。

水星是一颗“古典行星”，自古以来便被人类持续观测并识别为行星（或“游星”）。英文名称 Mercury 源自古罗马神祇 Mercurius，即商业与沟通之神，也是众神的信使。1974 年，“水手 10 号”（Mariner 10）完成了首次成功飞掠水星的任务，此后 MESSENGER 与 BepiColombo 轨道器也相继对其进行访问与探测。
\subsection{命名}

在历史上，人类会根据水星作为“昏星”或“晨星”出现的不同情况而称呼它不同的名字。到公元前约 350 年，古希腊人已经意识到这两颗“星”实际上是同一个天体[21]。他们称这颗行星为 Στίλβων（Stilbōn），意为“闪烁者”，以及 Ἑρμής（Hermēs），因其迅速而短暂的移动\,[22]；这一名称在现代希腊语中依然保留（Ερμής Ermis）[23]。罗马人则以他们的信使之神墨丘利（Mercury，拉丁语 Mercurius）为其命名，因为水星在天空中的移动速度比任何其他行星都快[21][24]。不过，根据老普林尼的记载，也有人将这颗行星与阿波罗相联系[25]。

水星的天文学符号是赫尔墨斯权杖（caduceus）的风格化形式；在 16 世纪，该符号加入了一个基督教十字，其形状为：☿[26][27]
\subsection{物理特性}
\begin{figure}[ht]
\centering
\includegraphics[width=8cm]{./figures/23fb842f59355847.png}
\caption{按比例绘制的水星及内太阳系具有行星质量的天体。从左至右依次为：水星、金星、地球、月球、火星和谷神星。} \label{fig_Mercur_2}
\end{figure}
水星是太阳系中四颗类地行星之一，这意味着它与地球一样，是一颗岩质天体。它是太阳系中最小的行星，赤道半径为 \(2{,}439.7\,\mathrm{km}\)（\(1{,}516.0\,\mathrm{mi}\)）\,[4]。水星的体积也小于——但质量大于——太阳系中最大的两个天然卫星：木卫三（盖尼米得）和土卫六（泰坦）。水星的组成大约为 70\% 的金属物质与 30\% 的硅酸盐物质\,[28]。
\subsubsection{内部结构}





